\Subhead{4.8.3}{Эффект Джозефсона}

Два сверхпроводника, слабо связанные через барьер, характеризуются разностью фаз $\phi$. В идеальном переходе
\begin{equation}
 I_s=I_c\sin\phi,\qquad \dot\phi=\frac{2eV}{\hbar}.
\end{equation}
Первое равенство описывает постоянный сверхток без напряжения; второе — изменение фазы при напряжении. При постоянном $V$ ток осциллирует с частотой $f_J=2eV/h$. Под внешним СВЧ-полем фазовая синхронизация даёт ступени напряжения $V_n=nhf/(2e)$ — эффект Шапиро и основа квантовых стандартов напряжения.

С учётом ёмкости и шунта модель RCSJ:
\begin{equation}
 I=I_c\sin\phi+\frac VR+C\dot V,
 \quad \frac{\hbar C}{2e}\ddot\phi+\frac{\hbar}{2eR}\dot\phi+I_c\sin\phi=I.
\end{equation}
Это аналог вынужденного затухающего маятника. При $I=0$ и малых фазах плазменная частота $\omega_J=\sqrt{2eI_c/(\hbar C)}$. Джозефсоновский генератор преобразует энергию источника напряжения в переменный ток и излучение; вывод мощности требует согласования с электромагнитной средой. Коэффициент $2e$ связан с зарядом пары.
